% Generated draft: requires mathematical and source review.
\documentclass[11pt]{article}
\usepackage[margin=25mm]{geometry}
\usepackage{fontspec}
\setmainfont{Latin Modern Roman}
\usepackage{amsmath,amssymb,mathtools,unicode-math}
\setmathfont{Latin Modern Math}
\usepackage{enumitem,needspace,xurl,hyperref}
\setlength{\parindent}{0pt}
\setlength{\parskip}{5pt}
\setlength{\emergencystretch}{3em}
\newcommand{\N}{\mathbb{N}}
\newcommand{\Z}{\mathbb{Z}}
\newcommand{\Q}{\mathbb{Q}}
\newcommand{\R}{\mathbb{R}}
\newcommand{\C}{\mathbb{C}}
\newcommand{\F}{\mathbb{F}}
\newcommand{\A}{\mathbb{A}}
\newcommand{\PP}{\mathbb{P}}
\newcommand{\E}{\mathbb{E}}
\newcommand{\eps}{\varepsilon}
\newcommand{\dd}{\,\mathrm{d}}
\begin{document}
\section*{UnsolvedMath ID: 800322}
\textbf{Source problem code:} catalogue entry
\subsection{Definitions and mathematical statement}
Let \((K,\partial)\) be a differential field of characteristic zero and let \(L\) be a differential field extension of \(K\). An element \(f\in L\) is \emph{differentially transcendental} (D-transcendental) over \(K\) if, for every integer \(r\geq 0\) and every nonzero polynomial \(P\in K[Y_0,\ldots,Y_r]\),
\[
P\bigl(f,\partial f,\ldots,\partial^r f\bigr)\neq 0.
\]
For functions or generating functions in a variable \(z\), the intended derivation is ordinarily differentiation with respect to \(z\); the coefficient field and the domain or formal-series setting must also be specified.

\textbf{Unresolved source statement.} The supplied record attributes Conjecture 6 to Pak (2018), but its mathematical display is truncated to \(\sum_{n\geq 1}\), followed by the words ``are D-transcendental.'' The summand, any additional generating functions, their variables, and the coefficient field are absent. Consequently, no particular generating function or complete conjectural assertion can be recovered from this record. The definition above supplies terminology only; it does not resolve these missing data.

The supplied background describes \(\Gamma(z)\) and \(\zeta(z)\) as known examples of D-transcendental functions and places the discussion in differential Galois theory. These are background claims in the record, not identified targets of the truncated conjecture. Its descriptions of the problems as open and unapproachable with existing tools are supplied source assessments, not independently verified current conclusions.
\subsection{Short English statement}
The record concerns a conjecture that certain generating functions are differentially transcendental, meaning that they satisfy no nonzero polynomial differential equation over a specified coefficient field. The formulas identifying those functions are missing, so the intended conjecture cannot yet be stated precisely. The supplied registry labels the problem open.
\subsection{Sources}
\begin{itemize}
\item Local UnsolvedMath registry, permanent ID \(800322\), source problem code \texttt{catalogue entry}, title ``Complexity in Enumerative Combinatorics --- Conjecture 6.'' The supplied record attributes the problem to Pak (2018) and labels it open.
\item Pak, \emph{Complexity Problems in Enumerative Combinatorics} (2018), Conjecture 6, \url{https://arxiv.org/abs/1803.06636}. The supplied provenance records access on 2026-10-16 and extraction by \texttt{arxiv\_latex}; neither access nor the extracted statement has been independently verified here.
\item The supplied background cites references labeled [Ritt] and [Put-S] for differential Galois theory and names Liouville, Lie, Picard, and Vessiot in its historical discussion. Full bibliographic details for these references are not supplied.
\end{itemize}
\end{document}
